% !PRINTED 579
\weberpage{579}
%<<CONT
the coefficients of $A(x)$, and is therefore an
integral rational function of $\lambda$ with rational numerical coefficients
of the form:
\weq{12}{S = S_0\lambda^s + S_1\lambda^{s-1} + S_2\lambda^{s-2} + \cdots + S_s,}
in which the $S_0, S_1, S_2, \ldots$, are rational numbers. That they are
\gesperrt{integers} is just what is to be proved. But this follows very simply if in
(12) one substitutes the development (3) for $\lambda$ and then compares
with one another the coefficients of equally high powers of $q$ on both
sides. For one obtains, for the determination of the $S_0, S_1, S_2, \ldots$,
a sequence of linear equations, each following one of which contains only
one new $S_i$, and that with the coefficient 1, while on the other side of
these equations there stand integral algebraic numbers. Accordingly the
$S_0, S_1, S_2, \ldots$, are integral algebraic numbers, and since they are
rational, they are also \gesperrt{integral rational numbers}.

Since we know moreover, by § 57, II and (7), $A(0)$ and $A(\infty)$,
there results for $A(x)$ the expression:
\weq{13}{A(x) = \pm x^{m^2-1} + a_1 x^{m^2-3} + a_2 x^{m^2-5} + \cdots + m,}
\gesperrt{in which the $a_1, a_2, \ldots$, are integral functions of
$\lambda$ with integral coefficients.}

Further, by § 57, (7):
\[
D(x) = \pm x^{m^2-1} A\left(\frac{1}{x}\right),
\]
and consequently:
\weq{14}{\pm D(x) = \pm 1 + a_1 x^2 + a_2 x^4 + \cdots + m x^{m^2-2}.}
\[
\text{[The sign } \pm \text{ in (13) and (14) is determined according to § 57 by}
(-1)^{\frac{m-1}{2}}\text{].}
\]
For the functions $B$ and $C$ the relation holds:
\weq{15}{B(x) = x^{m^2-1} C\left(\frac{1}{x}\right),}
and their roots are therefore reciprocal to one another. The first and
the last coefficient are equal to unity in $B$ and in $C$.

The roots of $B$ are:
\[
y_{h,h'} = \sqrt{\varkappa}\,\mathrm{sn}(\Omega_{h,h'}+K) = \sqrt{\varkappa}\,\frac{\mathrm{cn}\,\Omega_{h,h'}}{\mathrm{dn}\,\Omega_{h,h'}}
\]
\weq{16}{= \frac{\vartheta_{10}\left(\frac{h+h'\omega}{m},\omega\right)}{\vartheta_{00}\left(\frac{h+h'\omega}{m},\omega\right)}}
%CONT>>
